\documentclass[10pt,a4paper]{amsart}
\usepackage[margin=19mm]{geometry}
\usepackage{amsmath,amssymb,mathtools}
\usepackage[scr=rsfs,cal=euler]{mathalfa}
\usepackage{xfrac}
\usepackage[unicode,hidelinks,bookmarks=false]{hyperref}
\newcommand{\ie}{i.e.\ }
\newcommand{\cf}{cf.\ }
\newcommand{\resp}{respectively\ }
\newcommand{\Subsection}{\subsection}
\providecommand{\nobreakdash}{\nobreak\hbox{-}}
\setlength{\emergencystretch}{2em}
\multlinegap0pt
\usepackage{mathtools}
\usepackage{enumitem}
\usepackage{xurl}
\hypersetup{
    pdftitle={Hitting Arithmetic Progressions at the Square-Root Scale},
    pdfauthor={Samuel Korsky},
    pdfsubject={Arithmetic progressions; hitting sets},
    pdfkeywords={arithmetic progressions, hitting sets, progression-free sets}
}
\setlength{\parindent}{0pt}
\setlength{\parskip}{0.72em}
\newtheorem{theorem}{Theorem}[section]
\newtheorem{proposition}[theorem]{Proposition}
\newtheorem{lemma}[theorem]{Lemma}
\newtheorem{corollary}[theorem]{Corollary}
\newtheorem{conjecture}[theorem]{Conjecture}
\newtheorem{question}[theorem]{Question}
\newtheorem{remark}[theorem]{Remark}
\newcommand{\Z}{\mathbb Z}
\newcommand{\F}{\mathbb F}
\newcommand{\AP}{\mathrm{AP}}
\newcommand{\lcm}{\operatorname{lcm}}
\newcommand{\floor}[1]{\left\lfloor #1\right\rfloor}
\newcommand{\ceil}[1]{\left\lceil #1\right\rceil}
\numberwithin{equation}{section}
\begin{document}
\title[Hitting Arithmetic Progressions at the Square-Root Scale]{Hitting Arithmetic Progressions at the Square-Root Scale}
\author{Samuel Korsky}
\date{}
\maketitle
\noindent\emph{Source and licence.} Hitting Arithmetic Progressions at the Square-Root Scale. Version arXiv:2606.01780v1 (CC BY 4.0). This material is distributed under the Creative Commons Attribution 4.0 International licence, http://creativecommons.org/licenses/by/4.0/, as recorded in the arXiv OAI licence metadata for this exact version and shown on the abstract page https://arxiv.org/abs/2606.01780v1. Full rights notice: licenses/arxiv-2606.01780-CC-BY-4.0.txt.\par\smallskip
\noindent\textit{Editorial scope.} Counted unit: the original \texttt{theorem} environment \texttt{thm:below} at source lines 149--163 together with its complete proof at lines 622--738; the delivered document keeps source lines 84--739 so that every referenced label and every citation resolves inside it. Native editable source is the paper's own concatenated TeX; the AMS wrapper is editorial formatting only. No publisher class, upstream script or unrelated artwork is redistributed.
\section{Introduction}

For positive integers $N$ and $k$, let
\[
        [N]_0=\{0,1,\ldots,N-1\}.
\]
We write $f(N,k)$ for the minimum cardinality of a set $A\subseteq[N]_0$ which intersects every $k$-term arithmetic progression contained in $[N]_0$.
Throughout, arithmetic progressions are understood to have positive common difference.
The same function was introduced by Brown and Freedman \cite{BrownFreedman}, using the translated interval $\{1,\ldots,N\}$.
Equivalently, $[N]_0\setminus A$ contains no $k$-term arithmetic progression.

For fixed $k$, this problem is the complement of the classical extremal problem for progression-free sets.
If $r_k(N)$ denotes the largest size of a subset of $[N]_0$ with no $k$-term arithmetic progression, then
\[
        f(N,k)=N-r_k(N).
\]
Thus Szemer\'edi's theorem \cite{Szemeredi} says that $f(N,k)=N-o_k(N)$ for every fixed $k$.
Quantitative estimates for $r_k(N)$ form a central part of additive combinatorics, including Gowers's bounds for Szemer\'edi's theorem \cite{Gowers}, the recent Kelley--Meka and Bloom--Sisask progress on three-term progressions \cite{KelleyMeka,BloomSisask}, and the Leng--Sah--Sawhney bounds for longer progressions \cite{LengSahSawhney}.
Behrend--Rankin-type lower-bound constructions for $r_k(N)$ \cite{Behrend,Rankin}, including O'Bryant's formulation for long progressions \cite{OBryant}, also belong to this fixed-$k$ or slowly-growing-$k$ context.

The present paper is concerned with a different regime, where $k$ grows as a power of $N$.
Brown and Freedman proved, among other results, that if $k=N^\varepsilon$, then
\[
        f(N,k)\ll_\varepsilon N^{1-\varepsilon},
\]
matching the trivial lower bound $f(N,k)\ge \floor{N/k}$ up to a constant depending on $\varepsilon$.
They also studied the special prime-square case $N=p^2$, $k=p$, and proved
\[
        f(p^2,p)\le 2p-2
\]
for odd primes $p$.
Brown and Freedman also recorded Truss's improvement of the lower bound at the square-root scale.
Truss's published result states that
\[
        n+\frac12 n^{1/2}-2
        <
        f(n^2,n)
        \le
        p+\ceil{\frac{(n-1)^2}{p}},
\]
where $p$ is the largest prime at most $n$ \cite{Truss}.
The upper bound is asymptotically $2n+o(n)$.
Xu later obtained further upper-bound improvements for $r_p(p^2)$, equivalently upper-bound savings for $f(p^2,p)$, using a generalization of Szekeres's algorithm \cite{Xu}.

Our first result improves Truss's lower-bound constant.

\begin{theorem}[Critical lower bound]\label{thm:critical}
As $n\to\infty$,
\[
        f(n^2,n)
        \ge
        n+\left(\frac1{\sqrt2}+o(1)\right)n^{1/2}.
\]
\end{theorem}

The proof is elementary.
Partition $[n^2]_0$ into $n$ consecutive blocks of length $n$.
Every block is itself an $n$-term arithmetic progression, so every hitting set of size $n+s$ corresponds to at most $s$ nonsingleton blocks.
Thus there is a long run of consecutive blocks containing exactly one selected point each.
The relative positions of the selected points in such a run form a nonincreasing sequence, and its descent sequence satisfies strong divisibility constraints.
A finite congruence lemma shows that every nonconstant descent sequence has quadratic total mass.
If the descent sequence is constant, a separate residue-class argument gives the same asymptotic constant.

The same argument extends below the critical scale.


\begin{theorem}[Below-critical lower bound]\label{thm:below}
Fix $\theta\ge0$.
Let $k=k(n)\le n$ be a sequence of positive integers satisfying
\[
        n-k=\theta n^{1/2}+O(1).
\]
Then, as $n\to\infty$,
\[
        f(n^2,k)
        \ge
        \floor{\frac{n^2}{k}}
        +
        \left(\frac{\sqrt{\theta^2+2}-\theta}{2}+o(1)\right)n^{1/2}.
\]
\end{theorem}


Our upper-bound result improves the prime-square construction.

\begin{theorem}[Random front upper bound]\label{thm:upper}
For every sufficiently large prime $p$ we have
\[
        f(p^2,p)
        \le
        2p-\left(\sqrt{\frac23}-o(1)\right)\sqrt{\frac p{\log p}}.
\]
\end{theorem}

Here and throughout, $\log$ denotes the natural logarithm.
The construction is a randomized front construction followed by an alteration step.
One keeps the column-zero points only from row $h$ onward, places all large nonzero residues in row $h-1$, and randomly assigns the first $H\asymp h\log p$ residues among the first $h$ rows.
Large common differences are handled deterministically by the final front row.
For the remaining progressions, a deterministic reduction shows that they must be contained in the first $h$ rows and have common difference at most $h$; the expected number missed by the random part is then small enough to repair directly.

We also include the following orientation result, which explains why $k= N^{1/2}$ is a transition point.

\begin{proposition}[Above the square-root scale]\label{prop:above-root}
Fix $1/2<c<1$ and let $k=\floor{N^c}$.
Then
\[
        f(N,k)=(1+o(1))\cdot\frac Nk.
\]
\end{proposition}

The trivial block lower bound gives the lower estimate.
For the upper estimate, choose the largest prime $q\le k$.
By the prime number theorem,
\[
        q=(1+o(1))k.
\]
Since $c>1/2$, we have $(N-1)/(k-1)=o(k)$, so $q>(N-1)/(k-1)$ for all sufficiently large $N$.
Then take the multiples of $q$ in $[N]_0$.
This set has size $(1+o(1))N/k$.
It hits every $k$-term progression: if a progression has common difference $d\ge q$, then its span is at least $(k-1)q>N-1$, impossible; hence $1\le d<q$, so $d$ is invertible modulo $q$, and the first $q\le k$ terms already cover every residue class modulo $q$.
Thus the scale $k=N^{1/2}$ is the point at which the simple prime-modulus construction no longer automatically matches the block lower bound.

\section{A Congruence Lemma}

We first prove the arithmetic lemma used in the lower bound.

\begin{lemma}\label{lem:lcm}
For integers $a\ge1$ and $t\ge0$,
\[
        a\binom{a+t}{t}
        \mid
        \lcm(a,a+1,\ldots,a+t).
\]
\end{lemma}

\begin{proof}
Let
\[
        L=\lcm(a,a+1,\ldots,a+t).
\]
It is enough to compare $p$-adic valuations for each prime $p$.
For $j\ge1$, let $c_j$ be the number of integers in $\{a,a+1,\ldots,a+t\}$ divisible by $p^j$.
If $e=v_p(L)$, then $c_j=0$ for $j>e$.
For $1\le j\le e$,
\[
        c_j\le \floor{\frac{t}{p^j}}+1,
\]
because an interval of length $t+1$ contains at most $\floor{t/p^j}+1$ multiples of $p^j$.
Therefore
\[
\begin{aligned}
        v_p\!\left(\prod_{i=0}^t(a+i)\right)-v_p(t!)
        &=
        \sum_{j\ge1}c_j-
        \sum_{j\ge1}\floor{\frac{t}{p^j}} \\
        &\le
        \sum_{j=1}^e 1
        =e
        =v_p(L).
\end{aligned}
\]
Since
\[
        \frac{a(a+1)\cdots(a+t)}{t!}
        =
        a\binom{a+t}{t},
\]
the result follows.
\end{proof}

\begin{lemma}[Descent congruence lemma]\label{lem:sequence}
Let $m\ge1$, and let
\[
        \delta_0,\delta_1,\ldots,\delta_{m-1}
\]
be nonnegative integers.
Suppose that for every $1\le q\le m$ and every $0\le i\le m-q$,
\[
        q\mid \delta_i+\delta_{i+1}+\cdots+\delta_{i+q-1}.
\]
If $(\delta_i)$ is nonconstant, then
\[
        \sum_{i=0}^{m-1}\delta_i\ge \frac{m(m-1)}2.
\]
\end{lemma}

\begin{proof}
Subtract $\min_i\delta_i$ from every $\delta_i$.
This preserves the divisibility hypotheses, preserves nonconstancy, and can only decrease the sum.
Thus it suffices to prove the result under the additional assumption
\[
        \min_i\delta_i=0.
\]

Define prefix sums
\[
        P_0=0,
        \qquad
        P_j=\sum_{i=0}^{j-1}\delta_i
        \qquad (1\le j\le m).
\]
The hypothesis says that
\[
        j-i\mid P_j-P_i
        \qquad
        (0\le i<j\le m).
\]
In particular, $m\mid P_m$.
Write
\[
        P_m=am.
\]
We prove that if $(\delta_i)$ is nonconstant, then $a\ge (m-1)/2$.

Assume, for contradiction, that
\[
        a<\frac{m-1}{2}.
\]
Set
\[
        E_j=P_j-aj.
\]
Then
\[
        E_0=E_m=0
\]
and
\[
        j-i\mid E_j-E_i
        \qquad
        (0\le i<j\le m).
\]
Also
\[
        \delta_i=P_{i+1}-P_i=a+E_{i+1}-E_i\ge0.
\]

We prove by induction from the two endpoints inward that all $E_j$ vanish.
Suppose that for some integer $r\ge1$ with $2r<m$, we already know
\[
        E_0=\cdots=E_{r-1}=0
\]
and
\[
        E_{m-r+1}=\cdots=E_m=0.
\]
Put
\[
        h=m-2r>0.
\]
From divisibility by distances to the already-vanishing endpoint values, both $E_r$ and $E_{m-r}$ are divisible by
\[
        \Lambda=
        \lcm(1,2,\ldots,r,h+1,h+2,\ldots,h+r).
\]
Since
\[
        \Lambda\ge h+r=m-r>a,
\]
the inequality
\[
        \delta_{r-1}=a+E_r-E_{r-1}=a+E_r\ge0
\]
forces $E_r\ge0$.
Indeed, a negative multiple of $\Lambda$ would be at most $-\Lambda<-a$.
Similarly,
\[
        \delta_{m-r}=a+E_{m-r+1}-E_{m-r}=a-E_{m-r}\ge0
\]
forces $E_{m-r}\le0$.

If either $E_r$ or $E_{m-r}$ is nonzero, then
\[
        D:=E_r-E_{m-r}>0.
\]
The integer $D$ is divisible by $h=(m-r)-r$, and it is also divisible by $\Lambda$.
Therefore $D$ is divisible by $\lcm(h,\Lambda)$.
On the other hand,
\[
\begin{aligned}
        D
        &=
        \sum_{i=r}^{m-r-1}(E_i-E_{i+1}) \\
        &=
        \sum_{i=r}^{m-r-1}(a-\delta_i)
        \le ah.
\end{aligned}
\]
Thus
\[
        a\ge \frac{\lcm(h,\Lambda)}{h}.
\]
But $\lcm(h,\Lambda)$ is a multiple of $\lcm(h,h+1,\ldots,h+r)$.
By Lemma~\ref{lem:lcm},
\[
        \frac{\lcm(h,\Lambda)}{h}
        \ge
        \binom{h+r}{r}
        \ge
        h+r
        =m-r
        >\frac{m-1}{2},
\]
contradicting $a<(m-1)/2$.
Therefore
\[
        E_r=E_{m-r}=0.
\]
This completes the inward induction away from the possible middle point.

If $m$ is even, one middle value $E_{m/2}$ remains.
It is divisible by $\lcm(1,2,\ldots,m/2)$.
The inequalities on the two adjacent increments give
\[
        -a\le E_{m/2}\le a.
\]
Since $\lcm(1,2,\ldots,m/2)\ge m/2>a$, this forces $E_{m/2}=0$.
Thus $E_j=0$ for all $j$.

It follows that
\[
        \delta_i=P_{i+1}-P_i=a
\]
for all $i$.
Since $\min_i\delta_i=0$, we get $a=0$ and hence $\delta_i=0$ for all $i$, contradicting nonconstancy.
Therefore $a\ge(m-1)/2$, and
\[
        \sum_{i=0}^{m-1}\delta_i=P_m=am\ge \frac{m(m-1)}2.
\]
\end{proof}

\section{Singleton Runs}

Let $N=n^2$, and let $k\le n$.
Put
\[
        M=\floor{\frac{N}{k}}.
\]
Partition the initial interval $[0,Mk-1]$ into the $M$ consecutive blocks
\[
        I_j=[jk,(j+1)k-1],
        \qquad
        0\le j<M.
\]
Each $I_j$ is a $k$-term arithmetic progression with common difference $1$.

A block $I_j$ is called singleton if $|A\cap I_j|=1$.
A singleton run is a consecutive string of singleton blocks.

\begin{lemma}[Long singleton run]\label{lem:run}
Let $A\subseteq[N]_0$ meet every $k$-term arithmetic progression in $[N]_0$, and suppose
\[
        |A|=M+s.
\]
Then there is a singleton run of length
\[
        L\ge \frac{M-s}{s+1}.
\]
Writing the unique selected points in this run as
\[
        x_i=(u+i)k+r_i,
        \qquad
        0\le r_i<k,
        \qquad
        0\le i<L,
\]
and putting
\[
        m=L-1,
        \qquad
        \delta_i=k-(x_{i+1}-x_i)=r_i-r_{i+1}
        \quad (0\le i<m),
\]
one has
\[
        \delta_i\ge0,
        \qquad
        \sum_{i=0}^{m-1}\delta_i\le k-1,
\]
and
\[
        q\mid \delta_i+\delta_{i+1}+\cdots+\delta_{i+q-1}
\]
for every $1\le q\le m$ and every $0\le i\le m-q$.
\end{lemma}

\begin{proof}
Every block $I_j$ is a $k$-term progression, so every block contains at least one point of $A$.
Since $|A|=M+s$, at most $s$ blocks are nonsingleton.
Thus at least $M-s$ blocks are singleton, divided into at most $s+1$ runs.
One run has length at least $(M-s)/(s+1)$.

Now consider such a run.
If $x_{i+1}-x_i>k$, then the interval
\[
        [x_i+1,x_i+k]
\]
is a $k$-term arithmetic progression of common difference $1$ and contains no point of $A$.
This is impossible.
Therefore $x_{i+1}-x_i\le k$, so $\delta_i\ge0$.
The equality $\delta_i=r_i-r_{i+1}$ shows that the $r_i$ are nonincreasing, and hence
\[
        \sum_{i=0}^{m-1}\delta_i=r_0-r_m\le k-1.
\]

It remains to prove the divisibility condition.
Fix $q$ and $i$ with $1\le q\le m$ and $0\le i\le m-q$.
The union
\[
        I_{u+i}\cup I_{u+i+1}\cup\cdots\cup I_{u+i+q-1}
\]
is an interval of length $qk$.
For each residue class modulo $q$, the elements of this interval in that residue class form a $k$-term arithmetic progression with common difference $q$.
Since $A$ meets every such progression, the $q$ singleton points
\[
        x_i,x_{i+1},\ldots,x_{i+q-1}
\]
must occupy all residue classes modulo $q$.
Thus every $q$-window of selected singleton-run points is a complete residue system modulo $q$.
Comparing adjacent $q$-windows gives
\[
        x_{i+q}\equiv x_i\pmod q.
\]
But
\[
        x_{i+q}-x_i
        =qk-
        \sum_{j=i}^{i+q-1}\delta_j.
\]
Therefore
\[
        q\mid \sum_{j=i}^{i+q-1}\delta_j.
\]
\end{proof}

\section{Lower Bounds}

We first prove the critical result.

\begin{proof}[Proof of Theorem~\ref{thm:critical}]
Let $A\subseteq[n^2]_0$ meet every $n$-term arithmetic progression, and write
\[
        |A|=n+s.
\]
It suffices to prove
\[
        s\ge \left(\frac1{\sqrt2}+o(1)\right)n^{1/2}.
\]
If $s$ is not $O(n^{1/2})$ this is immediate, so assume $s=O(n^{1/2})$.

Apply Lemma~\ref{lem:run} with $k=n$ and $M=n$.
We obtain a singleton run of length $L=m+1$ with
\begin{equation}\label{eq:critical-run-length}
        m\ge \frac{n-s}{s+1}-1
        =
        \frac{n}{s+1}-O(1).
\end{equation}
Let $\delta_0,\ldots,\delta_{m-1}$ be the associated descent sequence.

First suppose this sequence is nonconstant.
By Lemmas~\ref{lem:sequence} and~\ref{lem:run},
\begin{equation}\label{eq:critical-nonconstant-descent}
        \frac{m(m-1)}2
        \le
        \sum_{i=0}^{m-1}\delta_i
        \le n-1.
\end{equation}
Thus
\begin{equation}\label{eq:critical-nonconstant-upper}
        m\le (\sqrt2+o(1))n^{1/2}.
\end{equation}
Combining \eqref{eq:critical-run-length} and \eqref{eq:critical-nonconstant-upper} gives
\[
        s\ge \left(\frac1{\sqrt2}+o(1)\right)n^{1/2}.
\]

It remains to handle the case where the descent sequence is constant.
Write
\[
        \delta_i=\delta
        \qquad
        (0\le i<m).
\]
Then
\[
        x_{i+1}-x_i=n-\delta=:d.
\]
All $m+1$ selected points in the singleton run are congruent modulo $d$.
Also
\begin{equation}\label{eq:critical-delta-bound}
        m\delta=\sum_i\delta_i\le n-1,
        \qquad
        \delta\le \frac{n-1}{m}.
\end{equation}

For every residue class $a$ modulo $d$, the progression
\[
        a,
        a+d,
        \ldots,
        a+(n-1)d
\]
lies in $[n^2]_0$, since $d\le n$.
Thus $A$ must meet every residue class modulo $d$.
The singleton run occupies only one residue class modulo $d$, so the points outside the run must cover at least $d-1$ other residue classes.
There are $n+s-(m+1)$ points outside the run, whence
\[
        n+s-(m+1)
        \ge
        d-1
        =n-\delta-1.
\]
Equivalently, using \eqref{eq:critical-delta-bound},
\begin{equation}\label{eq:critical-constant-bound}
        s\ge m-\delta
        \ge
        m-\frac{n-1}{m}.
\end{equation}
Combining \eqref{eq:critical-run-length} and \eqref{eq:critical-constant-bound} gives the desired bound.
Indeed, along any convergent subsequence write
\[
        s=c n^{1/2}+o(n^{1/2}).
\]
If $c=0$, then \eqref{eq:critical-run-length} gives $m/n^{1/2}\to\infty$, and \eqref{eq:critical-constant-bound} gives $s/n^{1/2}\to\infty$, a contradiction.
Thus $c>0$.
Now \eqref{eq:critical-run-length} gives
\[
        m\ge (1/c+o(1))n^{1/2}.
\]
If $c<1/\sqrt2$, then $m>\sqrt n$ for large $n$.
Since the function $t\mapsto t-(n-1)/t$ is increasing for $t>\sqrt{n-1}$, \eqref{eq:critical-constant-bound} gives
\[
        c\ge \frac1c-c,
\]
contradicting $c<1/\sqrt2$.
Thus every subsequential limit satisfies $c\ge1/\sqrt2$, and the theorem follows.
\end{proof}

We now prove the below-critical extension.


\begin{proof}[Proof of Theorem~\ref{thm:below}]
Let
\[
        N=n^2,
        \qquad
        M=\floor{\frac{N}{k}},
\]
and write
\[
        |A|=M+s.
\]
Since
\[
        k=n-\theta n^{1/2}+O(1),
\]
we have
\[
        M=n+\theta n^{1/2}+O(1),
\]
and hence
\begin{equation}\label{eq:below-delta-M}
        \Delta:=M-k=2\theta n^{1/2}+O(1).
\end{equation}

By Lemma~\ref{lem:run}, there is a singleton run of length $L=m+1$ with
\begin{equation}\label{eq:below-run-length}
        m\ge \frac{M-s}{s+1}-1
        =
        \frac{M}{s+1}-O(1).
\end{equation}
If $s$ is not $O(n^{1/2})$, the desired lower bound is immediate, so assume $s=O(n^{1/2})$.

If the associated descent sequence is nonconstant, then Lemmas~\ref{lem:sequence} and~\ref{lem:run} give
\begin{equation}\label{eq:below-nonconstant-descent}
        \frac{m(m-1)}2\le k-1,
\end{equation}
so
\begin{equation}\label{eq:below-nonconstant-upper}
        m\le (\sqrt2+o(1))n^{1/2}.
\end{equation}
Therefore, by \eqref{eq:below-run-length},
\[
        s\ge \left(\frac1{\sqrt2}+o(1)\right)n^{1/2}.
\]
Since
\[
        \frac{\sqrt{\theta^2+2}-\theta}{2}\le \frac1{\sqrt2},
\]
this is sufficient.

It remains to handle the constant-descent case.
Suppose
\[
        \delta_i=\delta
        \qquad
        (0\le i<m).
\]
Then
\[
        x_{i+1}-x_i=k-\delta=:d.
\]
Since $k\le n$, we have $d\le k\le n$.
For every residue class $a$ modulo $d$, the progression
\[
        a,
        a+d,
        \ldots,
        a+(k-1)d
\]
lies in $[n^2]_0$, because its largest possible last term is $kd-1\le k^2-1\le n^2-1$.
Thus $A$ must meet every residue class modulo $d$.
The singleton run occupies one residue class modulo $d$, so
\[
        M+s-(m+1)
        \ge
        d-1
        =k-\delta-1.
\]
Equivalently,
\[
        s\ge m+k-M-\delta
        =m-\Delta-\delta.
\]
Since $m\delta\le k-1$, we get
\begin{equation}\label{eq:below-constant-bound}
        s\ge m-\Delta-\frac{k-1}{m}.
\end{equation}

Let
\begin{equation}\label{eq:below-ctheta}
        c_\theta=\frac{\sqrt{\theta^2+2}-\theta}{2}.
\end{equation}
Combining \eqref{eq:below-run-length} and \eqref{eq:below-constant-bound} gives the claimed constant.
Indeed, along any convergent subsequence write
\[
        s=c n^{1/2}+o(n^{1/2}).
\]
If $c=0$, then \eqref{eq:below-run-length} gives $m/n^{1/2}\to\infty$, and \eqref{eq:below-constant-bound} gives $s/n^{1/2}\to\infty$, a contradiction.
Thus $c>0$.
Now \eqref{eq:below-run-length} gives
\[
        m\ge \left(\frac1c+o(1)\right)n^{1/2}.
\]
If $c<c_\theta$, then $c<1/\sqrt2$, so $m>\sqrt n$ for large $n$.
Since $t\mapsto t-(k-1)/t$ is increasing for $t>\sqrt{k-1}$, \eqref{eq:below-constant-bound} and \eqref{eq:below-delta-M} imply
\[
        c
        \ge
        \frac1c-2\theta-c.
\]
Equivalently,
\[
        2c^2+2\theta c-1\ge0.
\]
The positive root of this quadratic is $c_\theta$, contradicting $c<c_\theta$.
Thus every subsequential limit is at least $c_\theta$, and the theorem follows.
\end{proof}


\begin{thebibliography}{99}
\bibitem[Behrend]{Behrend} F. A. Behrend, \emph{On sequences of integers containing no arithmetic progression}, \v{C}asopis P\v{e}st. Mat. Fys. \textbf{67} (1938), 235--239. \href{https://eudml.org/doc/25161}{EuDML}.
\bibitem[BloomSisask]{BloomSisask} T. F. Bloom and O. Sisask, \emph{An improvement to the Kelley--Meka bounds on three-term arithmetic progressions}, \href{https://arxiv.org/abs/2309.02353}{arXiv:2309.02353}.
\bibitem[BrownFreedman]{BrownFreedman} T. C. Brown and A. R. Freedman, \emph{Small sets which meet all the $k(n)$-term arithmetic progressions in the interval $[1,n]$}, J. Combin. Theory Ser. A \textbf{51} (1989), 244--249. \href{https://doi.org/10.1016/0097-3165(89)90049-6}{doi:10.1016/0097-3165(89)90049-6}; \href{https://www.sfu.ca/~vjungic/tbrown/tom-38.pdf}{author PDF}.
\bibitem[Gowers]{Gowers} W. T. Gowers, \emph{A new proof of Szemer\'edi's theorem}, Geom. Funct. Anal. \textbf{11} (2001), 465--588. \href{https://doi.org/10.1007/s00039-001-0332-9}{doi:10.1007/s00039-001-0332-9}.
\bibitem[KelleyMeka]{KelleyMeka} Z. Kelley and R. Meka, \emph{Strong bounds for $3$-progressions}, Proceedings of the 64th IEEE Symposium on Foundations of Computer Science, 2023, 933--973; \href{https://arxiv.org/abs/2302.05537}{arXiv:2302.05537}.
\bibitem[LengSahSawhney]{LengSahSawhney} J. Leng, A. Sah, and M. Sawhney, \emph{Improved bounds for Szemer\'edi's theorem}, \href{https://arxiv.org/abs/2402.17995}{arXiv:2402.17995}.
\bibitem[OBryant]{OBryant} K. O'Bryant, \emph{Sets of integers that do not contain long arithmetic progressions}, Electron. J. Combin. \textbf{18} (2011), no. 1, Paper 59. \href{https://doi.org/10.37236/546}{doi:10.37236/546}; \href{https://arxiv.org/abs/0811.3057}{arXiv:0811.3057}.
\bibitem[Rankin]{Rankin} R. A. Rankin, \emph{Sets of integers containing not more than a given number of terms in arithmetical progression}, Proc. Roy. Soc. Edinburgh Sect. A \textbf{65} (1960), 332--344. \href{https://www.cambridge.org/core/journals/proceedings-of-the-royal-society-of-edinburgh-section-a-mathematics/article/xxivsets-of-integers-containing-not-more-than-a-given-number-of-terms-in-arithmetical-progression/7CB49707304932619FC2F3FB4A6BDE81}{Cambridge Core}.
\bibitem[Szemeredi]{Szemeredi} E. Szemer\'edi, \emph{On sets of integers containing no $k$ elements in arithmetic progression}, Acta Arith. \textbf{27} (1975), 199--245. \href{https://eudml.org/doc/205339}{EuDML}.
\bibitem[Truss]{Truss} J. K. Truss, \emph{Small sets which meet all the $n$-term arithmetic progressions in the interval $[1,n^2]$}, Bull. London Math. Soc. \textbf{23} (1991), 123--127. \href{https://doi.org/10.1112/blms/23.2.123}{doi:10.1112/blms/23.2.123}.
\bibitem[Xu]{Xu} X. Xu, \emph{A generalization of sets without long arithmetic progressions based on Szekeres algorithm}, J. Number Theory \textbf{133} (2013), 3670--3677. \href{https://doi.org/10.1016/j.jnt.2013.05.008}{doi:10.1016/j.jnt.2013.05.008}.
\end{thebibliography}
\end{document}
